\itemhead{2.2 例题 · 例题1（教材p.46）}{XJ-B1-C02-EX-P046-Y1243}
\label{XJ-B1-C02-EX-P046-Y1243}
{\footnotesize 湘教版教材p.46 / PDF第53页 / 2.2 例题 / 例题1}
\field{原题}{二次函数 $y=ax^2+bx+c$ 的图象如图 2.2-2 所示，顶点坐标为 $(2,2)$。根据图象解答下列问题：
\bankfig{C02-EX046}{0.38}
\subq{1}写出方程 $ax^2+bx+c=0$ 的两个根；
\subq{2}若方程 $ax^2+bx+c=k$ 有两个不相等的实数根，求 $k$ 的取值范围。}
\field{参考答案}{（1）由图得两个根为1、3；（2）抛物线开口向下且最大值2，与水平线 $y=k$ 有两个交点的条件为 $k<2$。}
{\small \field{关联知识}{Q01 二次函数、零点、判别式与根系关系（教材帮p.65 / PDF74）；F04 函数三种表示法、图象与读数（教材帮p.85、86、93、94 / PDF94、95、102、103）}
\field{方法与易错点}{按已知零点、顶点或系数选择交点式、顶点式或一般式；根平方和用根系关系消元。 易错：遗漏 $a
e0$；把二重根算成两个不同零点；有两正根只查积正。}
\field{答案依据}{例题原解位置：湘教版教材p.46及题库原图；独立编写可编辑参考解答。答案册此部分主要收练习、习题，不伪造例题答案册位置。}}
